\section{Unerwünschte Aussendungen}
\label{section:unerwuenschte_aussendungen_1}
\begin{frame}{Unerwünschte Aussendungen}
    \begin{center}
        \DARCimage{0.8\textwidth}{868include}
    \end{center}
\end{frame}

\begin{frame}{Unerwünschte Aussendungen}
    \begin{itemize}
        \item Amateurfunkverordnung (AFuV): \emph{Unerwünschte Aussendungen} auf das geringstmögliche Maß beschränken
        \item Es gibt weitere gesetzliche Regelungen für konkrete Grenzwerte
    \end{itemize}
\end{frame}


\begin{frame}
    \only<1>{
        \begin{QQuestion}{NJ201}{Ein Sender sollte so betrieben werden, dass~...}{er keine unerwünschten Aussendungen hervorruft.}
            {die Selbsterregung maximiert wird.}
            {parasitäre Schwingungen vorhanden sind.}
            {die Oberwellenabschirmung minimiert wird.}
        \end{QQuestion}

    }
    \only<2>{
        \begin{QQuestion}{NJ201}{Ein Sender sollte so betrieben werden, dass~...}{\textbf{\textcolor{DARCgreen}{er keine unerwünschten Aussendungen hervorruft.}}}
            {die Selbsterregung maximiert wird.}
            {parasitäre Schwingungen vorhanden sind.}
            {die Oberwellenabschirmung minimiert wird.}
        \end{QQuestion}

    }
\end{frame}

\begin{frame}
    \only<1>{
        \begin{QQuestion}{NF404}{Welche Eigenschaft sollte ein hinter einem VHF-Sender geschaltetes Filter haben? Dieses sollte...}{den gewünschten Frequenzbereich durchlassen.}
            {alle Oberschwingungen durchlassen.}
            {die Abstrahlung aller Nebenaussendungen zulassen.}
            {den gewünschten Frequenzbereich sperren.}
        \end{QQuestion}

    }
    \only<2>{
        \begin{QQuestion}{NF404}{Welche Eigenschaft sollte ein hinter einem VHF-Sender geschaltetes Filter haben? Dieses sollte...}{\textbf{\textcolor{DARCgreen}{den gewünschten Frequenzbereich durchlassen.}}}
            {alle Oberschwingungen durchlassen.}
            {die Abstrahlung aller Nebenaussendungen zulassen.}
            {den gewünschten Frequenzbereich sperren.}
        \end{QQuestion}

    }
\end{frame}

\begin{frame}
    \only<1>{
        \begin{QQuestion}{VD110}{Welche Aussage trifft die Amateurfunkverordnung (AFuV) hinsichtlich unerwünschter Aussendungen?}{Unerwünschte Aussendungen sind auf \qty{60}{\decibel} bezogen auf das Nutzsignal zu beschränken.}
            {Unerwünschte Aussendungen sind auf das geringstmögliche Maß zu beschränken.}
            {Unerwünschte Aussendungen sind auf \qty{40}{\decibel} bezogen auf das Nutzsignal zu beschränken.}
            {Unerwünschte Aussendungen sind nicht zulässig.}
        \end{QQuestion}

    }
    \only<2>{
        \begin{QQuestion}{VD110}{Welche Aussage trifft die Amateurfunkverordnung (AFuV) hinsichtlich unerwünschter Aussendungen?}{Unerwünschte Aussendungen sind auf \qty{60}{\decibel} bezogen auf das Nutzsignal zu beschränken.}
            {\textbf{\textcolor{DARCgreen}{Unerwünschte Aussendungen sind auf das geringstmögliche Maß zu beschränken.}}}
            {Unerwünschte Aussendungen sind auf \qty{40}{\decibel} bezogen auf das Nutzsignal zu beschränken.}
            {Unerwünschte Aussendungen sind nicht zulässig.}
        \end{QQuestion}

    }
\end{frame}%ENDCONTENT

\begin{frame}
    \frametitle{Messung von unerwünschten Aussendungen}
    \begin{center}
        \DARCimage{0.5\linewidth}{917include}
    \end{center}
    \begin{itemize}
        \item Am Senderausgang, unter Einzebiehung von Stehwellenmessgerät, Anpassgerät(e), Tiefpassfilter etc.
        \item Messsung von unerwünschen Aussendungen, die die Antenne erreichen können
    \end{itemize}
\end{frame}

\begin{frame}
    \only<1>{
        \begin{QQuestion}{EJ209}{Wie erfolgt die Messung der Leistungen, die zu unerwünschten Aussendungen führen?}{Die Messung erfolgt am Senderausgang mit einem hochohmigen HF-Tastkopf und angeschlossenem Transistorvoltmeter.}
            {Die Messung erfolgt am Fußpunkt der im Funkbetrieb verwendeten Antenne unter Einbeziehung des gegebenenfalls verwendeten Antennenanpassgeräts.}
            {Die Messung erfolgt am Ausgang der Antennenleitung unter Einbeziehung des im Funkbetrieb verwendeten Antennenanpassgeräts.}
            {Die Messung erfolgt am Senderausgang unter Einbeziehung des gegebenenfalls verwendeten Stehwellenmessgeräts und des gegebenenfalls verwendeten Tiefpassfilters.}
        \end{QQuestion}

    }
    \only<2>{
        \begin{QQuestion}{EJ209}{Wie erfolgt die Messung der Leistungen, die zu unerwünschten Aussendungen führen?}{Die Messung erfolgt am Senderausgang mit einem hochohmigen HF-Tastkopf und angeschlossenem Transistorvoltmeter.}
            {Die Messung erfolgt am Fußpunkt der im Funkbetrieb verwendeten Antenne unter Einbeziehung des gegebenenfalls verwendeten Antennenanpassgeräts.}
            {Die Messung erfolgt am Ausgang der Antennenleitung unter Einbeziehung des im Funkbetrieb verwendeten Antennenanpassgeräts.}
            {\textbf{\textcolor{DARCgreen}{Die Messung erfolgt am Senderausgang unter Einbeziehung des gegebenenfalls verwendeten Stehwellenmessgeräts und des gegebenenfalls verwendeten Tiefpassfilters.}}}
        \end{QQuestion}

    }
\end{frame}
